% !TeX root = ../../tesis.tex

% =============================================================================
%  SECCIÓN 3.3.10 — Clasificación Garibelt: instrumento de síntesis diagnóstica
% =============================================================================

\subsection{Clasificación Garibelt: instrumento de síntesis diagnóstica}
\label{subsec:clasificacion-garibelt}

Los ocho indicadores definidos en esta sección caracterizan la red de forma
independiente: cada uno mide una dimensión específica del grafo $\grafo$ y produce
un valor cuya unidad y escala le son propias. La Clasificación Garibelt opera en
un nivel distinto: no es el noveno indicador sino un instrumento de síntesis que
toma cinco de los ocho ---$C$, $C_i$, $DI$, $T$ y $B(v)$--- y los reúne en un
tablero de diagnóstico topológico. Su función dentro de esta sección es cerrar el
catálogo de indicadores con una herramienta que permita leerlos en conjunto sin
colapsar sus diferencias en un escalar único.

\begin{table}[htbp]
\centering
\caption[Clasificación Garibelt: dimensiones e instrumentos de normalización]{%
    Clasificación Garibelt v1: cinco dimensiones, indicadores fuente y métodos
    de normalización.}
\label{tab:clasificacion-garibelt-dims}
\begin{tabularx}{\textwidth}{c >{\raggedright\arraybackslash}X c c >{\raggedright\arraybackslash}X}
    \toprule
    \textbf{\#} & \textbf{Dimensión} & \textbf{Símbolo} &
    \textbf{Método de normalización} & \textbf{Referencia de calibración} \\
    \midrule
    1 & Accesibilidad — Cobertura Espacial
      & $C$
      & \textit{Literature Standard}
      & Estándar peatonal urbano (0–80\,\%) \\
    2 & Conectividad capilar — Fuerza Capilar
      & $C_i$
      & \textit{Empirical Percentile}
      & Percentiles Q1–Q3 del dataset \\
    3 & Eficiencia de ruta — Índice de Ruta Directa
      & $DI$
      & \textit{Theoretical Bounds}
      & $DI = 1$ óptimo; $DI \geq 2.5$ crítico \\
    4 & Fluidez global — Tiempo de Viaje Promedio
      & $T$
      & \textit{Reference-based}
      & \textsc{inegi}/\textsc{enoe} \cdmx{} (85–180\,min) \\
    5 & Centralidad crítica — Intermediación (Gini)
      & $B(v)$
      & $1 - \mathrm{Gini}(B)$
      & Monopolio topológico: Gini\,=\,0 equidistribuido \\
    \bottomrule
\end{tabularx}
\smallskip

\fuentefigura{Elaboración propia con base en \citet{bocarejo2012} y
\citet{talukder2017}.}
\end{table}

La Tabla~\ref{tab:clasificacion-garibelt-dims} asigna a cada dimensión su propio
método de normalización porque las cinco magnitudes pertenecen a naturalezas
estadísticas distintas. \citet{bocarejo2012} identifican cuatro componentes
inconmensurables en la accesibilidad urbana: infraestructura de transporte, uso
de suelo, restricciones temporales y condición del usuario; los cinco indicadores
del instrumento mapean directamente a los tres primeros. Llevarlos a una escala
común exige calibraciones distintas: referencia externa para $T$, límite teórico
para $DI$, percentil empírico para $C_i$ ---criterio que \citet{talukder2017}
avalan al señalar que la normalización debe ajustarse tanto al marco teórico como
a las propiedades del dato.

De los ocho indicadores del capítulo, dos quedan fuera de la versión~1 del
instrumento. $W_{s,h}$ (Penalización por Transferencia) tiene implementación
parcial en el motor \textit{VFTModel}; su incorporación queda diferida a la
versión~2. $\Delta E$ (Robustez Geométrica) requiere múltiples corridas del grafo
con nodos removidos, operación que corresponde a la Fase~4 y está fuera del
alcance del presente análisis. La aplicación del instrumento a la red de la
\zmvm{} ---incluyendo los valores normalizados y la lectura por dimensión--- se
presenta en la \autoref{sec:diagnostico-garibelt} del
\autoref{ch:analisis-red-actual}.
